% language=us runpath=texruns:manuals/ontarget

\pushoverloadmode

\input trac-riv.mkxl

\startcomponent ontarget-rivers

\environment ontarget-style

\startchapter[title={Crossing rivers}]

Occasionally you can hear or read a \TEX ie argue that the engine should be able
to deal with rivers. Now, to be honest, we never really bothered much about that,
if only because one should rewrite the content if one is really disturbed, but we
come to that. The excursion below is a side track of introducing some new graphic
capabilities in \CONTEXT\ while at the same time looking at possible improvements
in the paragraph builder.

Let's start with some terminology. In \type {tex.web} you can find Don Knuth
saying: \quotation {After considering \TEX's eyes and stomach, we come now to the
bowels} and I've heard talks at \TEX\ meetings where the presenter impresses the
audience by using these terms. The same happens with talking about widows, clubs,
penalties, demerits and such. Talking rivers fits into this phenomena. The fact that
after decades of \TEX\ usage no one really bothered much about it beyond mentioning
and suggesting half|-|way solutions is a telling sign.

If one talks about a river in a typeset text what is actually meant is that there
are spaces in successive lines that when you look at it seem to be connected and
that is considered bad or at least annoying. Of course there is nothing that
prevents us to consider (say) \type {m}'s or \type {oo}'s stacked in successive
lines to be seen as mountain rig but let's stick to spaces.

You can wonder what is so bad about a gently meandering river. In fact, the more
it meanders, the less it will be noticed. It's canals that matter more. In a
paragraph with wide spaces due to the constraints one can be annoyed by those,
and we will call them ponds. When they become too large we can talk about lakes.
Actually, when rivers are narrow, it would be better to talk of creeks.

So, in order to please those who love this analogy speak, we now have a decent
repertoire of seemingly bad things to look at: ponds and lakes, creeks and
rivers, and canals. An empty page is basically an ocean and too much white space
between paragraphs might be seen as a sea, but that can be intentional.

So how bad is all this in practice? Before we show some examples, we make a decision:
a pond is as bad as a creek, and a lake as bad as a river. Canals are even worse but
they will often show up like rivers: it is mostly a visual thing. A river can meander
which means that it can backtrack: a real worse case is a paragraph where it meanders
over multiple lines horizontally.

\startbuffer[profile]
01000000000000000000000000000001100000000000000
01000000000000000000011000000000011000000000011
01000001100000000001100110000000000110000001100
00100110011000000110000001100110000001100110000
00011000000110011000000000011001100000000110000
00000000000001100000000000000000000000000001100
00000000000000000000000000000000000000000110000
\stopbuffer

\startlocallinecorrection
\startMPcode
fill lmt_potraced [
    buffer  = "profile",
    size    = 1,
    value   = "0",
    polygon = true,
] xsized TextWidth
    withcolor "darkgreen" ;
\stopMPcode
\stoplocallinecorrection

Anyway, we distinguish three different cases small (ponds and creeks), large
(lakes and rivers) and whatever sits in between. In the next example we see these
three cases, where blue is small and red is large.

\startlocallinecorrection
\setbox\scratchbox\vbox{\samplefile{tufte}}%
\addriverstobox[criterium=7]\scratchbox
\box\scratchbox
\stoplocallinecorrection

The question now is, are these red rivers really that bad? Let's see how the text
looks when we don't identify them:

\samplefile{tufte}

We bet no one is really bothered by this. In fact, one can be more annoyed by the
commas at the end of nearly every line. This text (a quote from Tufte actually)
is also a bit atypical because of the many words separated by commas but that
makes it such an excellent quotation for many tests. There's always something a
user can complain about.

\startlocallinecorrection
\setbox\scratchbox\vbox{\samplefile{ward}}%
\addriverstobox[criterium=2]\scratchbox
\box\scratchbox
\stoplocallinecorrection

In the previous example we are in a way worse situation: plenty of creeks and
rivers there, but again, you won't hear us complain. So can we make an example
that will make us feel bad? In \in {figure} [fig:overflooding] we see something
that will never (or at least seldom) happen in a \TEX\ document: we see no
hyphenation.

\startplacefigure[title=Overflooding,reference=fig:overflooding]
\startcombination[nx=2,ny=1,location=top]
    {\setbox\scratchbox\vbox
       {\nohyphens
        \showglyphs
        \hsize 7cm
        \setupalign[verytolerant,stretch]%
        \samplefile{tufte}}%
      \addriverstobox[criterium=4]\scratchbox
      \box\scratchbox}
    {tufte}
    {\setbox\scratchbox\vbox
       {\nohyphens
        \showglyphs
        \hsize 7cm
        \setupalign[verytolerant,stretch]%
        \samplefile{knuth}}%
      \addriverstobox[criterium=7]\scratchbox
      \box\scratchbox}
    {knuth}
\stopcombination
\stopplacefigure

By now you might wonder how we get to these areas. What happens here is that we
first turn the typeset result into a (sort of) bitmap. Every line gets the same
number of slots, each representing 2pt advance. We start out with all slots being
zero but then assign slots that overlap with interword spacing to one. We can
then run over the rows and see where what columns have value one. But that is
not what we do. Because we want a picture, and because we also want meandering
visualized, we feed the bitmap into potrace, a library that is built in
\LUAMETATEX, and convert the resulting vector representation into a \METAPOST\
graphic.

If one is only interested if there are rivers at all, one can set the criteria
high and look only at the red blobs. In fact, if there is more than one red blob
one can already decide that we have a bad situation. But what can we do about it?
Well, not that much. Without hyphenation one is kind of lost, with hyphenation
one can best rewrite the paragraph and hope for the best. Of course a way out is
to use a new feature of \LUAMETATEX\ (and \CONTEXT): additional emergency par
builder steps. These come cheap and just try to do a better job by changing
constraints, one of which is selective expansion (aka hz). In fact, trying to get
rid of rivers and lakes by tweaking can be considered as harmful as messing with
looseness, a trick that one only applies when the results can be visually
checked. As with real rivers, messing with them will just make things worse.
Let's consider them natural boundaries that we don't even bother to cross.

\typebuffer[profile]

Just in case you're curious about what bitmap gets fed into potrace, we show it
above. You might see rivers, or not. And it might be clear now that we are not
going to waste time on this in the perspective of improving the par builder. It
would be one of these tuning options that in the end no one would use, just
because \quotation {fixing this} might as well \quotation {break that} which
means that such an automatism will demand closer inspection of the output than a
user wants. You might also realize that we don't really need potrace for
detecting lakes and rivers but it just gives nice graphics, and that's what we're
after. One can even argue that if avoiding any river is the objective, even one
pair of vertical ones is already enough to not look further and declare the whole
paragraph to be bad which then has to trigger a new attempt. It is however not
that hard to imagine something \typ {\riverdemerits} and maybe even \typ
{\doubleriverdemerits} based on a rather dumb fast pairwise check. The
complication of adding this to the par builder is that we then need to package
all the possible lines and that's not going to happen any time soon.

What might get unnoticed here is that we not only look at inter|-|word spacing
but also treat small glyphs like punctuation followed by such spacing as part of
a river and we can imagine plenty more cases to consider. It just shows that if
spacing is a problem, anything else can become one. We can also add some margin
to a space because being a bit more or less tolerant can matter.

% After wrapping up this experiment, Mikael Sundqvist found this example on stack
% exchange (entry 4507 52406) and after turning the example to use the same
% settings we get \in {figure} [fig:longrivers].
%
% \startbuffer
% % Of all the great rivers fo the world, none is as intriguing
% % as the Pearl. Short by world standards, it epitomizes the
% % old expression that good things come in small packages.
% Though the Pearl measures less than 50~miles in total length from its
% modest source as a cool mountain spring to the screaming cascades and
% steaming estuary of its downstream reaches, over those miles, the
% river has in one place or another everything you could possibly ask
% for. You can roam among lush temperate rain forests, turgid white
% water canyons, contemplative meanders among aisles of staid aspens
% (with trout leaping to slurp all the afternoon insects from its calm
% surface), and forbidding swamp land as formidable as any that Humphrey
% Bogart muddled through in \emph{The African Queen}.
% \stopbuffer
%
% \startplacefigure[title=A very long river,reference=fig:longrivers]
%     \scale [width=\textwidth] {%
%         \startcombination[nx=2,ny=1]
%           {\startshowrivers[criterium=8]
%              \switchtobodyfont[termes,10pt]
% %              \hangindent23pt \hangafter -3
%              \setupindenting[yes,23pt]
%              \showglyphs\hsize 247pt
%              \setupalign[verytolerant,stretch]
%              \getbuffer
%           \stopshowrivers} {shapes}
%           {\startshowrivers[criterium=8,option=test]
%              \switchtobodyfont[termes,10pt]
%              \hangindent23pt \hangafter -3
% %              \setupindenting[yes,23pt]
%              \showglyphs\hsize 247pt
%              \setupalign[verytolerant,stretch]
%              \getbuffer
%            \stopshowrivers} {polygons}
%         \stopcombination
%     }
% \stopplacefigure
%
% Indeed there is a pretty big canal running down the paragraph but is it really
% annoying enough to spend time on it? How many users stare long enough at a
% paragraph without reading it (which moves the eyes) to get fed up?

Just in case you want to experiment with this, here how it goes:

\starttyping
\startshowrivers[criterium=6]
    ...
\stopshowrivers
\stoptyping

where the following optional parameters can be set:

\starttabulate[]
\NC \type{step}           \NC \type{2pt}             \NC \NR
\NC \type{margin}         \NC \type{0pt}             \NC \NR
\NC \type{height}         \NC \type{.2\bodyfontsize} \NC \NR
\NC \type{depth}          \NC \type{.2\bodyfontsize} \NC \NR
\NC \type{spaceinbetween} \NC \type{.1\lineheight}   \NC \NR
\NC \type{criterium}      \NC \type{4}               \NC \NR
\stoptabulate

Here the \type {step} determines the accuracy, the \type {margin} is what gets
added to the space skip when it is encountered, the \type {height} and \type
{depth} are used to identify for instance punctuation characters that are treated
like glue. When we have a vertical glue that exceeds \typ {spaceinbetween} we
assume a new paragraph is started. Finally, \type {criterium} is what makes the
red shapes show up; steps 1 and 3 result in the green and blue ones.

Of course this is mostly an attempt to show that in practice we shouldn't care
too much about these rivers. As with quite some nice features, in the end once
they are there one forgets about them because in retrospect they were not that
much needed. This is why we only mention that there is a simple and fast variant
that one can enable with \typ {\showrivers} which acts upon the page but it is
less fun to look at.

Let us end by mentioning that we downloaded the Henry James book \quotation{The
turn of the screw} from the Project Gutenberg web page and compiled it in many
different page layouts and with different choices of fonts, and it was not so
easy to produce any disturbing rivers. \TEX\ does a very good job. The bigger
space after the period could be a bit disturbing, in particular if it appeared
on consecutive lines and if the lines had different stretch.

\startlocallinecorrection
\setbox\scratchbox\vbox{\samplefile{tufte}}
\addriverstobox[step=.5pt]\scratchbox
\ruledhbox{\box\scratchbox}
\stoplocallinecorrection

% \page
% \enabletrackers[typesetting.rivers]
% \dorecurse{10}{
%     \hsize \dimexpr\textwidth-#1cm\relax
%     \samplefile{tufte}\par
% }

% \page
% \enabletrackers[typesetting.rivers]
% \dorecurse{30}{
%     \samplefile{knuth}\par
%     \samplefile{ward}\par
%     \samplefile{tufte}\par
% }

\stopchapter

\stopcomponent
